\section{Extended Survey by Failure Mode and Intervention}
\label{sec:extended-survey}

This appendix expands the compact synthesis in the main paper. Its purpose is not to add a citation dump, but to explain how the broader evidence map is organized by failure mode, intervention point, and evaluation target. We group papers by their primary role in the framework; many papers naturally span multiple categories.

\subsection{Evidence-Centric Robustness: Noisy Retrieval and Distraction}

The first branch of robust RAG studies what enters the context. Early work on irrelevant context established that retrieval errors do not merely lower recall; they can actively redirect generation \citep{irrelevantContext2023}. Noise-aware training and robust RALM work then separated random negatives, hard negatives, counterfactual passages, and noisy retrieved evidence as different stressors rather than one generic perturbation \citep{raat2024,rrag037,rrag091}. Benchmark papers such as RGB and URAG made this decomposition operational by testing negative rejection, information integration, counterfactual evidence, and uncertainty under retrieval noise \citep{rgb2023,rrag001}. These works motivate the main-paper claim that ``noise robustness'' is not a single axis.

Distraction work focuses on negative marginal utility: a passage is harmful not simply because it is irrelevant, but because it changes the answer, confidence, attribution, or reasoning path. Long-context and context-utility studies show that more context can increase the probability of including plausible but misleading passages \citep{rrag003,rrag039,rrag032}. Utility-oriented retrieval and downstream-aware metrics therefore ask whether a passage helps the generator rather than whether it only matches the query \citep{rrag058,rrag067,rrag073,rrag081,rrag088}. Context compression and filtering methods try to remove irrelevant spans or documents while preserving evidence needed for multi-hop reasoning \citep{rrag044,chainOfNote2023,sufficientContext2024,rrag043,rrag048,rrag055}.

The later evidence-centric literature broadens the source of noise beyond the top-$k$ ranking. Query-level robustness work studies sensitivity to paraphrase and wording changes \citep{queryRobustness2025}; long-context RAG and context-position studies examine how ordering and context length alter evidence use \citep{rrag140,rrag143,rrag154,rrag166,rrag168,rrag170}. Other papers study context construction, compression, and retrieval defects as pipeline-level problems rather than isolated retriever errors \citep{rrag174,rrag179,rrag180,rrag181,rrag182,rrag183,rrag184}. More recent systems and benchmarks extend this line with stress tests for retrieval defects, context utility, and robust fine-tuning \citep{rrag200,rrag201,rrag206,rrag210,rrag212,rrag215,rrag218,rrag219,rrag224,rrag225,rrag226}. The common lesson is that retrieval should be evaluated as context construction.

\subsection{Model-Centric Evidence Use: Conflict, Faithfulness, and Attribution}

Once evidence enters the prompt, robustness depends on how the generator uses it. Knowledge-conflict papers study contradictions between parametric memory and retrieved evidence, or among retrieved sources \citep{tugOfWar2024,faithfulRAG2025,astuteRAG2025,rrag046}. These papers show that a model can over-trust retrieved evidence, resist correct retrieval because of parametric priors, or silently choose one source without exposing disagreement. Conflict-aware methods therefore localize discrepancies, compare internal and external knowledge, or introduce source-aware reasoning before generation \citep{rrag008,rrag013,rrag017,rrag024,rrag029,rrag031,rrag036,rrag038}.

Faithfulness and attribution papers study whether generated claims are supported by evidence. RAGTruth and related grounding work provide fine-grained hallucination or support labels, while source-preference and entity-conflict studies examine whether models cite and trust the right evidence \citep{ragTruth2024,rrag110,rrag112,rrag115,rrag132}. These papers connect robust RAG to citation verification: document-level citation is often too coarse, because a document can contain multiple claims, qualifiers, and contradictions. Claim-level or sentence-level attribution is more informative but increases annotation and verification cost \citep{rrag076,rrag077,rrag082,rrag086,rrag089,rrag092,rrag100}.

The broader conflict and faithfulness map includes work on source reliability, contextual faithfulness, real-document conflicts, and hidden-state or logit interventions \citep{rrag118,rrag130,rrag147,rrag150,rrag152,rrag163,rrag176,rrag187}. Newer papers extend conflict analysis to richer benchmark settings, source-aware generation, and practical RAG pipelines \citep{rrag202,rrag204,rrag207,rrag209,rrag214,rrag216,rrag220,rrag222,rrag223,rrag227}. Across these papers, the main-paper distinction between context availability and context use recurs: evidence can be present but ignored, cited but unsupported, or correct but less trusted than a stale parametric belief.

\subsection{Adversarial Evidence Corruption and Security}

Adversarial RAG work asks what happens when evidence unreliability is created by an adaptive attacker. Corpus and retrieval poisoning papers show that small numbers of malicious passages can manipulate dense retrieval or steer final answers \citep{poisoningRetrievalCorpora2023,poisonedRAG2024,badRAG2024,rrag049,rrag057}. Certified and aggregation-based defenses study robustness under bounded corruption assumptions \citep{certifiedRAG2024,rrag016,rrag020}, while reranking and detection defenses use gradient, attention, clustering, or anomaly signals to demote suspicious evidence \citep{rrag004,rrag009,rrag011,rrag015,rrag017}.

Prompt-injection and backdoor work targets a different interface. Instead of inserting false facts, the attacker may insert instructions that the model treats as commands, or triggers that activate malicious behavior through retrieved contexts \citep{aip2025,rrag005,trojanRAG2024,phantom2024,rrag010,rrag018}. These attacks require retrieval realism: a malicious document must be retrieved under plausible queries, not only under attacker-crafted prompts. This is why the main paper separates factual poisoning, indirect prompt injection, backdoor triggers, and prompt-template attacks as different threat games even when they all appear as retrieved text.

The security evidence map is broad and fast-moving. It includes poisoning benchmarks and arenas \citep{rrag019,rrag023,rrag033,rrag035}, traceback and forensics \citep{rrag014,rrag128}, privacy and membership inference \citep{rrag060,rrag063,rrag096,rrag103}, opinion or fairness manipulation \citep{rrag104,rrag162}, and graph, multimodal, or reasoning-style poisoning surfaces \citep{rrag051,rrag127,rrag175}. Many defenses and attacks are preprint-heavy \citep{rrag022,rrag025,rrag026,rrag047,rrag050,rrag052,rrag053,rrag054,rrag056,rrag059,rrag061,rrag062,rrag068,rrag070,rrag072,rrag078,rrag079,rrag080,rrag083,rrag084,rrag085,rrag090,rrag094,rrag097,rrag098,rrag099,rrag101,rrag102,rrag106,rrag107,rrag109,rrag114,rrag116,rrag117,rrag119,rrag120,rrag121,rrag122,rrag123,rrag124,rrag125,rrag126,rrag129,rrag133,rrag134,rrag137,rrag138,rrag139,rrag142,rrag144,rrag145,rrag148,rrag149,rrag155,rrag156,rrag157,rrag158,rrag160,rrag161,rrag165,rrag171,rrag172,rrag177,rrag185,rrag188,rrag189,rrag190,rrag194,rrag195,rrag196,rrag197,rrag199,rrag203,rrag208,rrag211,rrag213,rrag217,rrag221,rrag229}. The main use of this large set is not to rank attacks, but to motivate capability-normalized reporting.

\subsection{System Control: Uncertainty, Calibration, and Abstention}

System-control work converts evidence judgments into actions. Answerability and sufficient-context papers ask whether the retrieved context contains enough evidence to answer \citep{sufficientContext2024,rrag111,rrag141,rrag193}. Uncertainty and calibration work asks whether confidence reflects answer correctness, evidence sufficiency, retrieval usefulness, or selective risk \citep{calibRAG2025,uncertaintyRAG2024,rrag001,rrag007,rrag012,rrag093}. Adaptive retrieval work asks when the system should retrieve, retrieve again, or answer from internal knowledge \citep{selfRAG2023,whenToRetrieve2024,crag2024,ctrla2025,rrag034}.

The control literature highlights a recurring failure: confidence is not the same as evidence sufficiency. A model can be confident in an unsupported answer, or uncertain even when evidence is present but hard to integrate. Trust-or-abstain and self-aware RAG methods therefore model whether to answer, abstain, retrieve more, or ask for clarification \citep{rrag002,rrag006,rrag041,rrag064,rrag066,rrag074,rrag075,rrag087,rrag095,rrag105,rrag108}. Multi-turn and underspecified-query benchmarks add another layer: the correct action may be clarification before retrieval, not direct answering \citep{rrag141,rrag146,rrag153}.

Recent work extends control to cost-aware retrieval, conformal or set-valued prediction, risk-sensitive abstention, and adaptive policies over multiple retrieval steps \citep{rrag159,rrag164,rrag167,rrag169,rrag173,rrag178,rrag186,rrag191,rrag192,rrag198,rrag205,rrag228}. These papers motivate the main paper's action-control view: robust RAG is not only about reducing errors, but about choosing the right response under uncertainty and cost.

\subsection{Evaluation, Benchmarks, and Survey Context}

Evaluation papers provide the diagnostic infrastructure for robust RAG. RGB decomposes RAG reliability into multiple abilities \citep{rgb2023}; RAGTruth provides fine-grained hallucination annotations \citep{ragTruth2024}; query-robustness and context-utility work stress retrieval interfaces \citep{queryRobustness2025,rrag058}; security benchmarks and arenas evaluate poisoning and defense assumptions \citep{rrag019,rrag023,rrag035}; and uncertainty benchmarks evaluate whether systems know when their answers are reliable \citep{rrag001,rrag093}. Survey papers on RAG, trustworthy RAG, secure RAG, and evaluation provide the broader landscape \citep{ragSurvey2024,ragComprehensive2025,ragEvaluationSurvey2024,trustRAGCompass2024,trustworthyRAGSurvey2025,securityRAGSurvey2026,rrag134,rrag151}.

What is still missing is joint evaluation. Existing benchmarks often isolate one failure mode: noise, conflict, poisoning, or unanswerability. Real systems face combinations. A benchmark that includes a stale source, a hard negative, insufficient evidence, and a malicious document will test a different capability from any single-axis benchmark. This motivates the action-trace blueprint: robust-RAG evaluation should score evidence state, answer correctness, action choice, cost, and threat assumptions together.

